\subsection{LOWER CENTRAL SERIES, NILPOTENT GROUPS}

Let \(\mathbf{G}\) be a group, \(\mathbf{H}\) a subgroup of \(\mathbf{G}\) and \(\mathbf{K}\) a normal subgroup of \(\mathbf{G}\) . The image of \(\mathbf{H}\) in \(\mathbf{G} / \mathbf{K}\) is contained in the centre of \(\mathbf{G} / \mathbf{K}\) if and only if \((\mathbf{G}, \mathbf{H}) \subset \mathbf{K}\) .

DEFINITION 5. Let G be a group. The lower central series of G is the sequence \((\mathbf{C}^n (\mathbf{G}))_{n\geqslant 1}\) of subgroups of G defined inductively by:

\[
\mathbf {C} ^ {1} (\mathbf {G}) = \mathbf {G}, \quad \mathbf {C} ^ {n + 1} (\mathbf {G}) = (\mathbf {G}, \mathbf {C} ^ {n} (\mathbf {G})).
\]

Let \(f: \mathbf{G} \to \mathbf{G}'\) be a group homomorphism. It is seen, by induction on \(n\) , that \(f(\mathbf{C}^n(\mathbf{G})) \subset \mathbf{C}^n(\mathbf{G}')\) and that, if \(f\) is surjective, \(f(\mathbf{C}^n(\mathbf{G})) = \mathbf{C}^n(\mathbf{G}')\) . In particular, for all \(n \geqslant 1\) , \(\mathbf{C}^2(\mathbf{G})\) is a characteristic (and hence normal) subgroup of \(\mathbf{G}\) . For all \(n \geqslant 1\) , \(\mathbf{C}^n(\mathbf{G}) / \mathbf{C}^{n+1}(\mathbf{G})\) is contained in the centre of \(\mathbf{G} / \mathbf{C}^{n+1}(\mathbf{G})\) .

Let \((\mathbf{G}_1, \mathbf{G}_2, \ldots)\) be a decreasing sequence of normal subgroups of \(\mathbf{G}\) such that (1) \(\mathbf{G}_1 = \mathbf{G}\) ; (2) for all \(i\) , \(\mathbf{G}_i / \mathbf{G}_{i+1}\) is contained in the centre of \(\mathbf{G} / \mathbf{G}_{i+1}\) . Then \(\mathbf{C}^i(\mathbf{G}) \subset \mathbf{G}_i\) , as is seen by induction on \(i\) .

Now

\[
\left(\mathbf {C} ^ {m} (\mathbf {G}), \mathbf {C} ^ {n} (\mathbf {G})\right) \subset \mathbf {C} ^ {m + n} (\mathbf {G}).\tag{7}
\]

For, if this relation is denoted by \((\mathbf{F}_{m,n})\) , it follows from \((\mathbf{F}_{m,n})\) , by no. 2, Proposition 5, that

\[
\begin{array}{l} \left(\mathrm{C} ^ {m} (\mathrm{G}), \mathrm{C} ^ {n + 1} (\mathrm{G})\right) \subset \left(\mathrm{G}, \left(\mathrm{C} ^ {m} (\mathrm{G}), \mathrm{C} ^ {n} (\mathrm{G})\right)\right). \left(\mathrm{C} ^ {n} (\mathrm{G}), \left(\mathrm{G}, \mathrm{C} ^ {m} (\mathrm{G})\right)\right) \\ \subset \mathrm{C} ^ {m + n + 1} (\mathrm{G}). \left(\mathrm{C} ^ {m + 1} (\mathrm{G}), \mathrm{C} ^ {n} (\mathrm{G})\right). \end{array}
\]

Hence \(((\mathbf{F}_{m,n})\) and \((\mathbf{F}_{m + 1,n}))\Rightarrow (\mathbf{F}_{m,n + 1})\) . As \((\mathbf{F}_{m,1})\) and \((\mathbf{F}_{1,n})\) are obvious \((\mathbf{F}_{m,n})\) follows by induction.

DEFINITION 6. A group G is called nilpotent if there exists an integer n such that \(\mathbf{C}^{n+1}(\mathbf{G}) = \{e\}\) . The least integer n such that \(\mathbf{C}^{n+1}(\mathbf{G}) = \{e\}\) is called the nilpotency class of a nilpotent group G.

If \(n \in N\) , a group of nilpotency class n is called a nilpotent group of class n.~It is sometimes said that the nilpotency class of a group G is finite if G is nilpotent.

Examples. (1) A group is nilpotent of class 0 (resp. \(\leqslant 1\) ) if and only if it consists of the identity element (resp. is commutative).

\begin{enumerate}
\def\labelenumi{(\arabic{enumi})}
\setcounter{enumi}{1}
\item
  *For every commutative ring A and every integer \(n \geqslant 1\) , the upper strict triangular group \(\mathrm{T}_1(n, \mathbf{A})\) is nilpotent of class \(\leqslant n - 1\) (and exactly of class \(n - 1\) if \(\mathbf{A} \neq \{0\}\) ).*
\item
  Let \(\mathbf{G}\) be a nilpotent group of class \(n\) . Every subgroup (resp. every quotient group) of \(\mathbf{G}\) is nilpotent of class \(\leqslant n\) . For, if \(\mathbf{H}\) is a subgroup of \(\mathbf{G}\) , then \(\mathbf{C}^n (\mathbf{H})\subset \mathbf{C}^n (\mathbf{G})\) . If \(\mathbf{G}'\) is a quotient group of \(\mathbf{G}\) and \(\pi : \mathbf{G}\to \mathbf{G}'\) is the canonical homomorphism, then \(\mathbf{C}^n (\mathbf{G}') = \pi (\mathbf{C}^n (\mathbf{G}))\) .
\item
  A finite product of nilpotent groups is nilpotent.
\end{enumerate}

PROPOSITION 7. Let G be a group and n an integer. The following conditions are equivalent:

\begin{enumerate}
\def\labelenumi{(\alph{enumi})}
\item
  G is nilpotent of class \(\leqslant n\) .
\item
  There exists a series of subgroups of \(\mathbf{G}\) :
\end{enumerate}

\[
\mathrm{G} = \mathrm{G} ^ {1} \supset \mathrm{G} ^ {2} \supset \dots \supset \mathrm{G} ^ {n + 1} = \{e \}
\]

such that \((\mathbf{G},\mathbf{G}^k)\subset \mathbf{G}^{k + 1}\) for all \(k\in [1,n]\) .

\begin{enumerate}
\def\labelenumi{(\alph{enumi})}
\setcounter{enumi}{2}
\item
  There exists a subgroup \(\mathbf{A}\) of \(\mathbf{G}\) contained in the centre of \(\mathbf{G}\) such that \(\mathbf{G} / \mathbf{A}\) is nilpotent of class \(\leqslant n - 1\) .
\item
  (a) \(\Rightarrow\) (b): it suffices to take \(\mathbf{G}^k = \mathbf{C}^k (\mathbf{G})\) .
\item
  (b) \(\Rightarrow\) (a): by induction on \(k\) , \(\mathbf{C}^k (\mathbf{G})\subset \mathbf{G}^k\)
\item
  (a) \(\Rightarrow\) (c): it suffices to take \(\mathbf{A} = \mathbf{C}^n (\mathbf{G})\) .
\item
  (c) \(\Rightarrow\) (a): let \(\pi: \mathbf{G} \to \mathbf{G} / \mathbf{A}\) be the canonical homomorphism; then \(\pi(\mathbf{C}^n(\mathbf{G})) = \mathbf{C}^n(\mathbf{G} / \mathbf{A}) = \{e\}\) and hence \(\mathbf{C}^n(\mathbf{G}) \subset \mathbf{A}\) , whence \(\mathbf{C}^{n+1}(\mathbf{G}) = \{e\}\) .
\end{enumerate}

More briefly: a group is nilpotent of class \(\leqslant n\) if it can be obtained from the group \(\{e\}\) by n successive central extensions.

COROLLARY. A central extension of a nilpotent group (by a necessarily commutative group) is nilpotent.

PROPOSITION 8. Let G be a nilpotent group of class \(\leqslant n\) and let H be a subgroup of G. There exists a sequence of subgroups

\[
\mathbf {G} = \mathbf {H} ^ {1} \supset \mathbf {H} ^ {2} \supset \dots \supset \mathbf {H} ^ {n + 1} = \mathbf {H},
\]

such that \(\mathbf{H}^{k + 1}\) is normal in \(\mathbf{H}^k\) and \(\mathbf{H}^k /\mathbf{H}^{k + 1}\) is commutative for all \(k\leqslant n\)

Choose a sequence \((\mathbf{G}^{k})\) of subgroups of G satisfying the conditions of Proposition 7 (b) for all k; \(G^{k}\) is normal in G. Write:

\[
\mathbf {H} ^ {k} = \mathbf {H}. \mathbf {G} ^ {k}.
\]

It is necessary to verify that \(\mathbf{H}^{k + 1}\) is normalized by \(\mathbf{H}^k = \mathbf{H}.\mathbf{G}^k\) ; as it is normalized by \(\mathbf{H}\) , it suffices to verify that it is by \(\mathbf{G}^k\) . Now, if \(s\in \mathbf{G}^k\) and \(h\in \mathbf{H}\) ,

\[
s h s ^ {- 1} = s h s ^ {- 1} h ^ {- 1}. h \in (\mathrm{G}, \mathrm{G} ^ {k}). \mathrm{H}
\]

and \((\mathbf{G},\mathbf{G}^k)\) . H is contained in \(\mathbf{G}^{k + 1}\) . \(\mathbf{H} = \mathbf{H}^{k + 1}\) ; hence \(s.\mathbf{H}^{k + 1}.s^{-1} = \mathbf{H}^{k + 1}\) , which shows that \(\mathbf{H}^{k + 1}\) is normal in \(\mathbf{H}^k\) .

Finally, the canonical homomorphism \(G^{k}/G^{k+1} \to H^{k}/H^{k+1}\) is obviously surjective; as the first group is commutative, so is the second.

COROLLARY 1. Let G be a nilpotent group and H a subgroup of G. If H is distinct from G, the normalizer \(N_{G}(H)\) of H in G is distinct from H.

Let \(k\) be the largest index such that \(H^{k} \neq H\) . The group \(H^{k}\) normalizes \(H\) and is distinct from \(H\) .

COROLLARY 2. Let G be a nilpotent group and H a subgroup of G. If H is distinct from G, there exists a normal subgroup N of G, containing H, distinct from G and such that G/N is commutative.

Let \(k\) be the least index such that \(H^k \neq G\) . The group \(H^k\) satisfies the required conditions.

COROLLARY 3. Let G be a nilpotent group and H a subgroup of G. If G = H.(G, G), then G = H.

Every subgroup N of G which contains H and such that G/N is commutative contains H.(G, G). Corollary 3 thus follows from Corollary 2.

Corollary 3 can also be formulated thus: let X be a subset of G. For X to generate G, it is necessary and sufficient that the image of X in G/D(G) generate G/D(G).

COROLLARY 4. Let \(f: \mathbf{G}' \to \mathbf{G}\) be a group homomorphism. Suppose that

\begin{enumerate}
\def\labelenumi{(\alph{enumi})}
\item
  G is nilpotent.
\item
  The homomorphism \(f_1 \colon \mathbf{G}' / (\mathbf{G}', \mathbf{G}') \to \mathbf{G} / (\mathbf{G}, \mathbf{G})\) , derived from \(f\) by passing to the quotients, is surjective.
\end{enumerate}

Then \(f\) is surjective.

This follows from Corollary 3 applied to the subgroup \(\mathbf{H} = f(\mathbf{G}')\) .

PROPOSITION 9. Let G be a nilpotent group of class \(\leqslant n\) and let N be a normal subgroup of G. There exists a series of subgroups

\[
\mathrm{N} = \mathrm{N} ^ {1} \supset \mathrm{N} ^ {2} \supset \dots \supset \mathrm{N} ^ {n + 1} = \{e \}
\]

such that \((\mathbf{G},\mathbf{N}^k)\subset \mathbf{N}^{k + 1}\) for \(k = 1,\dots ,n\)

If \((\mathbf{G}^k)\) satisfies condition (b) of Proposition 7, then take

\[
\mathbf {N} ^ {k} = \mathbf {G} ^ {k} \cap \mathbf {N}.
\]

COROLLARY 1. Let G be a nilpotent group, Z the centre of G and N a normal subgroup of G. If N ≠ \{e\}, then N ∩ Z ≠ \{e\}.

Let \(k\) be the largest index such that \(\mathbf{N}^k \neq \{e\}\) . The group \(\mathbf{N}^k\) is contained in \(\mathbf{N}\) . On the other hand, \((\mathbf{G}, \mathbf{N}^k) \subset \mathbf{N}^{k+1} = \{e\}\) ; hence \(\mathbf{N}^k\) is contained in the centre \(\mathbf{Z}\) of \(\mathbf{G}\) .

COROLLARY 2. Let \(f\) be a homomorphism from a nilpotent group \(G\) to a group \(G'\) . If the restriction of \(f\) to the centre of \(G\) is injective, \(f\) is injective.

This is Corollary 1 applied to \(\operatorname{Ker}(f)\) .

